% Generated draft: requires mathematical and source review.
\documentclass[11pt]{article}
\usepackage[margin=25mm]{geometry}
\usepackage{fontspec}
\setmainfont{Latin Modern Roman}
\usepackage{amsmath,amssymb,mathtools,unicode-math}
\setmathfont{Latin Modern Math}
\usepackage{enumitem,needspace,xurl,hyperref}
\setlength{\parindent}{0pt}
\setlength{\parskip}{5pt}
\setlength{\emergencystretch}{3em}
\newcommand{\N}{\mathbb{N}}
\newcommand{\Z}{\mathbb{Z}}
\newcommand{\Q}{\mathbb{Q}}
\newcommand{\R}{\mathbb{R}}
\newcommand{\C}{\mathbb{C}}
\newcommand{\F}{\mathbb{F}}
\newcommand{\A}{\mathbb{A}}
\newcommand{\PP}{\mathbb{P}}
\newcommand{\E}{\mathbb{E}}
\newcommand{\eps}{\varepsilon}
\newcommand{\dd}{\,\mathrm{d}}
\begin{document}
\section*{UnsolvedMath ID: 20002838}
\textbf{Source problem code:} catalogue entry
\subsection{Definitions and mathematical statement}
The assigned Problem 3 asks for cohomological tools for pseudorepresentation deformation rings analogous to those for Galois representations. The following Section 9 on Kedlaya's $(\varphi,\Gamma)$-modules and later Emerton material are extraction spillover, not additional hypotheses or questions here.

The aid describes a modern answer in a multiplicity-free residual setting. Fix a profinite group $\Gamma$, a coefficient field $k$ of characteristic $p$, and two distinct continuous characters $\chi_1,\chi_2:\Gamma\to k^\times$. Let the residual degree-two determinant pseudorepresentation be that of $\bar\rho=\chi_1\oplus\chi_2$. Use continuous cohomology and the finite-cohomology/representability hypotheses of the Cayley--Hamilton deformation setup; these hypotheses are not fully listed in the source. Write
\[
 \alpha=\chi_1\chi_2^{-1},\qquad
 B=H^1(\Gamma,k(\alpha)),\qquad
 C=H^1(\Gamma,k(\alpha^{-1})).
\]
A Cayley--Hamilton generalized matrix algebra retains the two off-diagonal extension modules and their compositions, whereas trace and determinant record closed paths rather than isolated off-diagonal arrows.

The aid's first-order visibility calculation is explicit. For a cocycle $b$ representing $B$, the upper triangular deformation over $k[\varepsilon]/(\varepsilon^2)$ is
\[
 \rho_b(g)=\begin{pmatrix}
  \chi_1(g)&\varepsilon b(g)\chi_2(g)\\0&\chi_2(g)
 \end{pmatrix}.
\]
Its trace and determinant equal the residual ones. The analogous lower triangular extension from $C$ is likewise invisible. More generally, in
\[
 (I+\varepsilon\begin{pmatrix}a(g)&b(g)\\c(g)&d(g)\end{pmatrix})
       \operatorname{diag}(\chi_1(g),\chi_2(g)),
\]
only the diagonal entries contribute to trace and determinant at first order. The first off-diagonal invariant involves the closed-cycle tensor $b\otimes c$; it has representation-parameter complexity two, even though an invariant cycle can give a first-order pseudodeformation coordinate. It is not the nonzero product of two $\varepsilon$ entries in the displayed dual-number matrix.

The two ordered diagonal compatibility classes are
\[
 B\otimes_k C\longrightarrow H^2(\Gamma,k)\oplus H^2(\Gamma,k),
 \qquad b\otimes c\longmapsto(b\smile c,\ c\smile b).
\]
Under the stated multiplicity-free deformation hypotheses, the aid identifies the complexity-two pseudodeformation tangent quotient, after removing split-character tangent directions, with the kernel of this pair of cup maps. The order and coefficient twists are retained; no independence of the two cohomology components is assumed. If $H^2(\Gamma,k)=0$ and $B,C$ are finite dimensional, its dimension is exactly $\dim_kB\,\dim_kC$. This is a tangent-dimension statement, not a claim that every higher-order pseudodeformation is unobstructed.

More generally the reported framework uses cup products and higher $A_\infty$ operations on extension spaces to describe relations and lifting obstructions. The record does not supply a complete general theorem for repeated residual constituents, arbitrary coefficient characteristics or every local/global deformation condition. In particular a determinant-law definition, rather than an unqualified trace-only convention, is required where those notions differ. The result supplies a scoped cohomological dictionary and the two-character visibility test, rather than a complete answer for all pseudorepresentation rings.
\subsection{Short English statement}
Cayley--Hamilton extension data provide cohomological tools near multiplicity-free pseudorepresentations. For two distinct residual characters, separate upper and lower extensions are invisible to first-order trace and determinant; closed cycles first appear through a tensor product constrained by two ordered cup classes.
\subsection{Sources}
\begin{itemize}
\item Michael Volpato's notes for the AIM workshop \emph{$p$-adic representations, modularity, and beyond}, Section 8, ``Kisin: Pseudo-representations,'' Problem 3, \url{https://aimath.org/WWN/padicmodularity/padicmodularity.pdf}.
\item Local UnsolvedMath registry, permanent ID 20002838, source code \texttt{catalogue entry}, canonical location \texttt{aim-representation-theory-notes.json}, \texttt{notes[109]}, record dated 14 August 2026. The statement is marked reconstructed unverified and clean statement null; verification identifies the short cohomological request before unrelated Section 9 and later-section spillover. The original literature field is empty.
\item Supplied research attempt 1: partial result, problem status at run partially solved, rigor labelled proved, novelty confidence low. Its generalized-matrix/cup-product dictionary and visibility-order lemma are machine-generated research aid, not peer reviewed; independently verify before citation. No complete modern bibliographic reference or full representability hypothesis list is supplied.
\end{itemize}
\end{document}
